\section{Soil structure and its formation}
\kw{Soil structure} is the arrangement of the primary particles into aggregates. It forms when the particles are brought together (by wetting and drying, freezing and thawing, the pressure of roots, the burrowing of worms) and are held by \emph{binding agents}: clay and oxides, the polysaccharides produced by microbes and roots, the hyphae of fungi, and humus. The stability of an aggregate to the impact of rain and to rapid wetting is its \kw{aggregate stability}; a stable soil keeps its pores, and an unstable one slakes, and the pores are sealed by a \kw{crust} that cuts the infiltration.
\par\noindent \textbf{Notes:} 1. Organic matter is the main agent that stabilises the aggregates of the topsoil; each loss of 1 per cent of organic carbon lowers the stability noticeably in a silt loam. 2. The large pores between aggregates carry the fast flow; the small pores inside the aggregates store water for the crop. A good structure has both.

\section{Tillage systems}
\kw{Tillage} is the mechanical working of the soil. The aims are to prepare a seedbed, to bury the residues, to control weeds and to loosen compacted soil. The systems are classified by the disturbance.
\begin{center}\small\begin{tabular}{@{}p{2.8cm}p{4.1cm}p{3.1cm}p{3.5cm}@{}}\toprule
\hd{System} & \hd{Operation} & \hd{Residue cover} & \hd{Effect on water}\\\midrule
Conventional & plough, then harrow; full inversion & less than 15 per cent & rough surface stores water at first, then crusts; fast drying of the seedbed\\
Reduced (minimum) & chisel or disc; partial & 15 to 30 per cent & intermediate\\
No-till & seed drilled into the stubble; no soil movement & more than 30 per cent & high infiltration through continuous pores; less evaporation; slower warming\\\bottomrule
\end{tabular}\end{center}
\noindent \textbf{Advantages of reduced tillage:} less erosion, less evaporation, savings in fuel and time, more organic matter near the surface. \textbf{Disadvantages:} more compaction near the surface in wet soils, a cooler seedbed in spring, a need for herbicide, and sometimes a yield lower in the first years. The results of the ensemble in this chapter show how the tillage changes the pore structure and the drainage.

\section{Compaction}
\kw{Compaction} is the increase of the bulk density by the load of machines or animals, with the loss of the large pores. The sign is a bulk density above a critical value for the texture (about 1.4 to 1.5 g\,cm$^{-3}$ in a clay loam, and 1.7 in a sand), at which the roots cannot penetrate. The consequences are a lower infiltration, a poor aeration, a shallow root system, more runoff and waterlogging. The remedy is prevention (less traffic on wet soil, controlled traffic lanes) because loosening by a subsoiler is short-lived.

\section{Pedotransfer functions}
The measurement of the retention curve and of the conductivity is slow and costly. A \kw{pedotransfer function} (PTF) estimates the hydraulic parameters from properties that are easy to find: the percentages of sand, silt and clay, the organic carbon and the bulk density. Examples are the regressions of Saxton and Rawls (Chapter~\ref{chap:11}) and the neural networks of Rosetta. A PTF gives the average behaviour of soils of similar texture; the error of the predicted conductivity is often one order of magnitude, so that it is suitable for planning, not for the final design of a field.

\section{Mulch: kinds, effects, advantages and disadvantages}
A \kw{mulch} is a layer of material placed on the soil surface. The kinds are \emph{organic} (straw, crop residues, wood chips, grass cuttings), \emph{inorganic} (plastic film, gravel) and \emph{living} (a cover crop). A mulch changes the water balance by: (1) cutting the energy that reaches the soil, so that the evaporation falls; (2) intercepting some rain (a loose mulch of 5 to 10 t\,ha$^{-1}$ holds 2 to 3 mm); (3) slowing the runoff and the impact of drops, and so protecting the structure; (4) changing the temperature of the soil (organic mulch cools it, a black film warms it).
\par\noindent \textbf{Advantages:} less evaporation, more water for the crop, control of weeds, less erosion, a more even soil temperature, more organic matter if it is organic. \textbf{Disadvantages:} cost and labour; a cold soil in spring; a habitat for slugs and rodents; the disposal of plastic; and, as Chapter~\ref{chap:09} shows, the water that is saved does not always become yield (it can drain).

\section{Evaporation from bare soil: the two-stage model}
Evaporation from a wet bare soil takes place in two stages (Ritchie, 1972). In \kw{stage 1} (the energy-limited stage) the surface is wet and the rate is the potential rate set by the energy; it lasts until a cumulative amount $U$ (6 to 12 mm, larger in a clay) has evaporated. In \kw{stage 2} (the soil-limited or falling-rate stage) the surface is dry, the vapour must diffuse through a dry layer, and the cumulative evaporation grows as the square root of the time:
\begin{equation*}\boxed{E(t)=U+\beta\sqrt{t-t_1}}\end{equation*}
with $\beta$ a soil constant (mm\,d$^{-1/2}$) and $t_1$ the time at which stage 1 ends. A small rain re-starts stage 1 and, because a small depth is lost quickly, the evaporation after small showers is a large share of the rain.

\section{Salinity and sodicity}
\kw{Salinity} is the content of soluble salts, measured by the electrical conductivity of the saturation extract $EC_e$ (dS\,m$^{-1}$). Crops are classified by the threshold above which the yield falls: tolerant (barley, cotton, $>6$), moderate (maize, wheat, 1.7 to 6), sensitive (bean, carrot, $<1.7$). \kw{Sodicity} is the excess of sodium on the exchange sites of the clay, measured by the sodium adsorption ratio $SAR=\mathrm{Na}/\sqrt{(\mathrm{Ca}+\mathrm{Mg})/2}$ (concentrations in meq\,L$^{-1}$); a high SAR disperses the clay, closes the pores and lowers the conductivity. A soil is \emph{saline} when $EC_e>4$ dS\,m$^{-1}$ and \emph{sodic} when $SAR>13$. Salinity is cured by leaching, sodicity by adding calcium (gypsum) before leaching.
